\chapter{Анализ сферы подбора персонала, методов обработки естественного языка и подходов семантического сопоставления вакансий и резюме}
\label{chapter:analysis}
\section{Системный анализ предметной области подбора персонала, процессов подбора персонала и систем автоматизированного подбора персонала}

Являясь программным ядром современного подбора персонала, системы автоматизированного подбора персонала поддерживают как минимум два критически важных и взаимосвязанных процесса: управление воронкой подбора персонала (от формирования длинного списка кандидатов к короткому списку) и поддержку принятия решений посредством ранжирования списков кандидатов и фильтрации больших наборов данных о соискателях \cite{src02}.

С точки зрения системного анализа, подбор персонала можно представить как потоковую систему обработки входящих заявок. Входом системы являются цифровые резюме соискателей и связанные с ними цифровые профили, включающие структурированные атрибуты опыта работы, профессиональных навыков, уровня образования, зарплатных ожиданий и статуса прохождения воронки подбора персонала, а выходом будут являться решения о продвижении соискателя по стадиям воронки и формирование на ранних этапах воронки упорядоченного списка соискателей для изучения рекрутером, чтобы он мог выбрать подходящих на должность кандидатов. Особенность данной предметной области заключается в том, что ранжирование списков соискателей относительно требований вакансий представляет собой высокорисковое управленческое решение (high-stakes decision): в отличие от ранжирования электронных документов, соискатели обладают собственными правами и интересами, а алгоритмическое ранжирование выступает инструментом поддержки принятия решений человеком, а не автономной системой выбора/отбора соискателей \cite{src02}.

Практический эффект алгоритмов ранжирования на популяцию кандидатов объясняется преимущественно моделью просмотра списков: позиции, находящиеся в топе выдачи, получают непропорционально большую долю внимания (экспозицию). В контексте управления персоналом данный эффект способен многократно усиливать различия между группами соискателей даже при незначительных отличиях в исходных оценках релевантности \cite{src02}. Таким образом, справедливость систем автоматизированного подбора персонала зависит не только от лежащего в их основе алгоритма, но и от интерфейса самой системы автоматизированного подбора персонала, а также от стратегии просмотра списков пользователем, которым может выступать например сотрудник по подбору персонала \cite{src02}.

Исторически инструменты, ориентированные на интеграцию с системами автоматизированного подбора персонала и автоматизацию процессов отбора кандидатов на работу, чаще всего использовали методы извлечения структурированных атрибутов из цифровых профилей соискателей (профессиональных навыков, стажа работы, уровня образования), механизмы фасетного поиска и техники ранжирования из арсенала классического информационного поиска. Характерным примером подобного подхода может служить система PROSPECT, позиционируемая как система поддержки принятия решений. Данный комплекс извлекает ключевые аспекты профиля кандидата, представляет их в виде фасетов и применяет алгоритмы информационного поиска для ранжирования списков соискателей относительно требований вакансии. В экспериментальной части исследования авторы сравнивали ранжирование с использованием извлеченной структурированной информации и вариант без такого обогащения профиля; по опубликованному описанию IBM Research добавление извлеченных аспектов профиля кандидата дало примерно 30\% улучшения качества ранжирования в задаче отбора кандидатов при найме \cite{src03}.

Современная эволюция систем автоматизированного подбора персонала характеризуется переходом к семантическому сопоставлению цифровых текстовых профилей соискателей и текстовых описаний вакансий на базе плотных векторных представлений, то есть числовых описаний смысла текста, и к обучению ранжирующих моделей на основе поведенческой обратной связи (кликов пользователей в системе, конверсий на следующие этапы воронки подбора персонала, прочих событий в системе автоматизированного подбора персонала, в том числе фактов принятия на работу конкретного соискателя в конкретную компанию). Современные индустриальные решения используют конвейеры многозадачного обучения моделей, а также алгоритмы на базе больших языковых моделей для более глубокого анализа требований вакансии и профилей соискателей. В частности, передовые архитектуры программного обеспечения поддержки процесса подбора персонала включают трехступенчатый процесс поиска талантов (отбор кандидатов в список, предварительное ранжирование списков кандидатов, окончательное переранжирование списков кандидатов), применяют большие языковые модели для извлечения предпочтений рекрутера из текстовых описаний вакансий и истории закрытия вакансий (успешного приема на работу сотрудника на позицию), а также задействуют модели класса смешения экспертов (mixture-of-experts) с учетом специфики конкретной роли/позиции/должности, на которую осуществляется поиск работника \cite{src17}.

В контексте настоящего исследования целевая система автоматизированного подбора персонала под названием Reqcore рассматривается как программная среда, для работы внутри которой модуль семантического сопоставления текстов вакансии и резюме должен выполнять следующие функции:

\begin{compactenum}
  \item Принимать на вход текстовое описание вакансии и набор текстов цифровых профилей кандидатов.
  \item Вычислять меру семантической близости между требованиями вакансии и содержанием резюме с учётом особенностей предметной области подбора персонала.
  \item Формировать ранжированный список кандидатов для первичного отбора на открытую вакансию.
  \item Обеспечивать объяснимость результатов выдачи соискателей в списках, например через сопоставление извлеченных из цифровых профилей кандидатов навыков и атрибутов.
  \item Предоставлять верифицируемую оценку качества работы алгоритмов ранжирования списков соискателей \cite{src02, src23}.
\end{compactenum}

\section{Аналитический обзор методов обработки естественного языка (Natural Language Processing) и методов машинного обучения для построения ранжирующей функции (Learning to Rank)}

Для решения задач сопоставления текстов вакансий и резюме традиционно применяются различные семейства методов представления и сравнения текстовой информации: лексические модели, такие как модель частоты термина с обратной частотой документа (Term Frequency-Inverse Document Frequency, TF-IDF), вероятностные ранжирующие функции, включая вероятностную функцию ранжирования документов под названием BM25, а в последнее время используються нейросетевые плотные векторные представления на базе моделей с механизмом самовнимания. В классической вероятностной парадигме функция ранжирования документов под названием BM25 зарекомендовала себя как надежный контрольный ориентир информационного поиска. Связь показателя релевантности с вероятностными предпосылками, параметрами насыщения частоты термина (term frequency, TF) и нормализацией по длине документа критически важна для текстов вакансий и резюме, поскольку они зачастую существенно различаются по объему и стилистике \cite{src11}. В современных системах поиска информации и рекомендаций, в том числе построенных на методологии обучения ранжированию (Learning to Rank), вероятностная функция ранжирования документов под названием BM25 по-прежнему служит устойчивым ориентиром; при интеграции нейросетевых алгоритмов ранжирования в контур системы автоматизированного подбора персонала требуется воспроизводимый классический контрольный метод для сравнительного анализа, и в данной работе таким методом выступает BM25 \cite{src08, src11}.

Методология обучения ранжированию (Learning to Rank) формализует процесс построения ранжирующей функции по заранее размеченным данным: уровням релевантности, парам предпочтения или спискам объектов с известным желаемым порядком. Общепринято классифицировать алгоритмы обучения ранжированию на три основные группы: поточечный подход (pointwise), сводящийся к регрессии или классификации по отдельным объектам, попарный подход (pairwise), обучающийся предсказывать предпочтения между двумя документами, и списочный подход (listwise), оптимизирующий функцию потерь непосредственно на всем списке документов на основе ранжирующих метрик \cite{src08}. Классическим примером попарного подхода выступает метод попарного обучения ранжированию под названием RankNet, в рамках которого задается вероятностная модель предпочтения между парой объектов исходя из разности их индивидуальных оценок, после чего минимизируется кросс-энтропийная ошибка функции потерь логистического связывания. Главное практическое преимущество подобных методов заключается в их дифференцируемости и встроенной совместимости с нейросетевыми архитектурами \cite{src09}.

Ключевым аспектом обучения моделей ранжирования является оценка качества получаемого ранжирования. Такие метрики, как дисконтированная кумулятивная выгода (Discounted Cumulative Gain, DCG) и нормализованная дисконтированная кумулятивная выгода (Normalized Discounted Cumulative Gain, NDCG), не только учитывают градуированный уровень релевантности, но и применяют прогрессивный дисконт в зависимости от занятой позиции, что полностью согласуется с когнитивным эффектом позиционного смещения при просмотре списков пользователем \cite{src10}. В прикладной сфере подбора персонала метрики семейства NDCG признаются наиболее подходящими, так как специалисты по подбору персонала редко просматривают весь перечень соискателей целиком: релевантность первых позиций ранжирования имеет прямое влияние на итоговый результат набора \cite{src02, src16}.

Самостоятельным фундаментальным направлением исследований алгоритмов машинного обучения выступает обеспечение справедливости ранжирования (fairness in ranking). В задачах такого типа список рассматривается не только как упорядоченный набор оценок, но и как механизм распределения ограниченного ресурса видимости: чем выше кандидат расположен, тем выше вероятность, что его профиль будет прочитан и отобран. Поэтому в моделях справедливого ранжирования вводятся ограничения на экспозицию относительно профессионально обоснованных признаков качества кандидата, например совпадения требуемых навыков, релевантного опыта, уровня позиции, образования или сертификаций. Метрика полезности в этом контексте обозначает целевую метрику качества выдачи, например нормализованную дисконтированную кумулятивную выгоду (Normalized Discounted Cumulative Gain, NDCG), а критериями справедливости могут быть ограничение дисбаланса экспозиции между группами, пропорциональность видимости квалификации кандидатов или отсутствие систематического занижения позиций защищенной группы \cite{src16, src36, src37}. Для архитектуры системы автоматизированного подбора персонала это представляет значительную концептуальную ценность: вероятность найма кандидата напрямую зависит от занимаемой позиции, а базовые числовые оценки соответствия, то есть исходные баллы, по которым алгоритм сортирует кандидатов перед показом пользователю, могут резко изменить распределение экспозиции даже при небольших различиях \cite{src02, src16}.

\section{Обоснование выбора архитектуры двунаправленных кодировочных представлений (Bidirectional Encoder Representations from Transformers), модели векторных представлений предложений (Sentence Embeddings using Siamese Bidirectional Encoder Representations from Transformers Networks) и методов векторизации}

Современным стандартом анализа семантики текста в области обработки естественного языка выступают модели с механизмом самовнимания, позволяющим эффективно моделировать смысловые зависимости внутри последовательностей знаков (tokens) без использования классических сверточных или рекуррентных подходов \cite{src07}. На основе этой парадигмы архитектура двунаправленных кодировочных представлений (Bidirectional Encoder Representations from Transformers) предложила схему двунаправленного предобучения языковых представлений на больших объемах неразмеченных текстов с их последующей тонкой настройкой под прикладные задачи (fine-tuning) \cite{src05}.

При решении задач семантического сопоставления текстов вакансии и резюме важнейшим архитектурным решением является выбор между перекрестным кодировщиком представлений (cross-encoder) и архитектурой раздельного кодирования текстов (bi-encoder). Перекрестный кодировщик производит совместную контекстную обработку пары текстов сквозь всю сеть, требуя высоких вычислительных затрат на формирование предсказания. Архитектура раздельного кодирования сначала преобразует текст вакансии и текст резюме в отдельные плотные векторные представления, то есть числовые векторы смысла, после чего сравнивает эти векторы с помощью стандартной косинусной меры сходства или скалярного произведения. В основополагающей статье по модели векторных представлений предложений (Sentence Embeddings using Siamese Bidirectional Encoder Representations from Transformers Networks) подробно раскрыто ограничение перекрестных кодировщиков: под большой базой понимается коллекция из тысяч или миллионов текстов, где попарное сравнение каждого запроса с каждым документом требует отдельного запуска модели для каждой пары. Если сравнивать 10 000 предложений друг с другом попарно, требуется около 50 млн сравнений; именно такой порядок вычислений делает перекрестный кодировщик непрактичным для интерактивного поиска по крупному набору резюме \cite{src06}. Модели семейства Sentence-BERT преобразуют архитектуру двунаправленного кодирования в сиамскую или триплетную архитектуру для генерации оптимизированных векторных представлений предложений, сравнимых напрямую; качество проверяется на наборах семантического сходства и перефразирования, например STS Benchmark и SICK-R, по корреляции Спирмена или Пирсона между модельной оценкой и экспертной разметкой \cite{src06, src29}.

В сценарии разработки модуля ранжирования цифровых профилей соискателей для системы автоматизированного подбора персонала вышеупомянутые выводы делают рациональной реализацию архитектуры двухступенчатого ранжирования. Первая ступень это стадия первичного поиска и формирования набора кандидатов, основанная на поиске по плотным векторным представлениям через архитектуру раздельного кодирования (bi-encoder) для быстрого отбора первых N кандидатов из общей базы цифровых профилей, то есть из всего множества резюме, доступных системе для данной организации или источника данных. Вторая ступень это стадия повторного ранжирования, которая может осуществляться более ресурсоемкой моделью перекрестного кодировщика либо строиться как отдельная надстройка обучения ранжированию (Learning to Rank) над широким набором признаков. Такая каскадная схема поиска соответствует индустриальной логике начинающийся с первичной фильтрация кандидатов затем предварительное ранжирование и наконец финальное ранжирование: сначала формируется ограниченный набор потенциально подходящих соискателей, затем он сортируется более точной моделью. В исследовании корпоративного поиска талантов LinkedIn такая схема описана как практический способ работать с большим числом профилей: сначала система отбирает управляемое число кандидатов из крупной базы, а затем применяет более дорогие признаки и модели для уточнения порядка выдачи \cite{src38}. Современная работа по ранжированию талантов также использует многоступенчатую схему, где крупная языковая модель помогает детализировать требования вакансии, а специализированная модель ранжирования уточняет порядок кандидатов по роли \cite{src17}.

Эмпирическую применимость решения на базе модели векторных представлений предложений (Sentence Embeddings using Siamese Bidirectional Encoder Representations from Transformers Networks) к предметной области найма доказывает не отдельная базовая архитектура, а доменно адаптированный подход CareerBERT: Matching Resumes to ESCO Jobs in a Shared Embedding Space for Generic Job Recommendations. Авторы переводят резюме соискателей и описания профессий в единое векторное пространство, где близкие по смыслу тексты располагаются рядом, а несвязанные тексты оказываются дальше. Важно, что обучающие пары в CareerBERT формируются иначе, чем в настоящей работе: авторы используют не исторические решения специалистов по откликам, а стандартизированную базу ESCO и создают пары из названия профессии и навыка, название профессии и синонима, и название профессии и описание. В статье указаны соответственно 113 685, 14 723 и 2 937 таких пар. После этого модель оценивается в два шага: на объявлениях EURES и на реальных резюме с обратной связью экспертов по управлению персоналом. Это показывает, что профессиональная таксономия может служить источником обучающих пар, когда прямой истории решений по откликам нет, тогда как в данной работе аналогичную роль выполняют события системы автоматизированного подбора персонала \cite{src01, src30}. Родственные подходы к сопоставлению соискателя и вакансии также используют глубокие сиамские сети, пары текст вакансии текст резюме, размеченные консультантами по подбору персонала, и подбор сложных отрицательных примеров, то есть похожих, но фактически неподходящих пар, чтобы модель лучше различала близкие профессиональные профили \cite{src04, src39, src40}.

При работе с русскоязычными текстами дополнительными факторами выступают качество готовых весовых моделей для целевого языка, а также доступность стандартных контрольных наборов для внешней проверки качества. Такой контрольный набор нужен не для обучения модели, а для воспроизводимого сравнения: он задает одинаковые задачи и метрики, позволяющие проверить разные модели в сопоставимых условиях до внедрения в прикладной контур. В современных публикациях вводится многоязычная методика массовой оценки плотных текстовых представлений (Massive Text Embedding Benchmark, MTEB), а для русского языка соответственно обширный контрольный набор оценки плотных текстовых представлений ruMTEB; в таких наборах формализуются задачи классификации, кластеризации, поиска, переранжирования и семантического сходства, что важно для выбора базовой модели и последующей проверки качества разрабатываемого модуля ранжирования соискателей \cite{src19, src20, src29}.

Выбор методов векторизации обязательно должен учитывать риски алгоритмического смещения: аудиты в системах подбора персонала и обзоры справедливости больших языковых моделей подтверждают, что использование только векторных плотных представлений текстов может провоцировать дискриминационные искажения результатов (по полу, расе, национальности и объему текстового цифрового профиля кандидата) \cite{src21, src22}.

\section{Выводы по первому разделу}

\begin{compactenum}
  \item Предметная область подбора персонала требует ранжирования кандидатов не как автономного кадрового решения, а как инструмента поддержки специалиста, поскольку итоговая позиция кандидата влияет на видимость профиля и дальнейшее рассмотрение.
  \item Классические методы информационного поиска, включая TF-IDF и BM25, остаются важным контрольным ориентиром, но недостаточно устойчивы к синонимии и вариативности профессионального языка в вакансиях и резюме.
  \item Архитектура раздельного кодирования текстов рациональна для прототипа, потому что позволяет заранее получать векторные представления вакансий и резюме и сравнивать их значительно дешевле, чем перекрестный кодировщик.
  \item Для оценки качества выбранного подхода необходимо использовать метрики ранжирования, учитывающие позиции кандидатов в списке, прежде всего NDCG, Precision@K, Recall@K и MRR.
\end{compactenum}

\section{Постановка задачи}

На основании проведенного анализа необходимо разработать программный модуль семантического сопоставления текстов вакансии и резюме, который принимает подготовленное описание вакансии и набор цифровых профилей кандидатов, строит плотные векторные представления текстов, рассчитывает числовую оценку их смысловой близости и возвращает упорядоченный список кандидатов для первичного рассмотрения специалистом по подбору персонала.

Входными данными задачи являются тексты вакансий, тексты резюме и исторические события подбора персонала, преобразованные в обучающие пары с бинарной меткой кадрового решения. Выходными данными являются ранжированный список кандидатов, числовые оценки соответствия и воспроизводимые экспериментальные метрики качества ранжирования на отложенной тестовой выборке.

Разрабатываемое решение должно удовлетворять трем ограничениям: сохранять контроль кадрового решения за человеком, быть воспроизводимым по исходным данным и программным артефактам, а также обеспечивать приемлемый компромисс между качеством ранжирования и стоимостью вывода по сравнению с более дорогими попарными стратегиями сопоставления.
